\begin{question}{30}{
    Bij militaire transportvluchten heeft het gewicht van de lading een directe invloed rol op het brandstofverbruik van een toestel.
    Naarmate er meer materieel of voorraden worden meegenomen, stijgt het verbruik tijdens de cruise-fase van de vlucht aanzienlijk.
    In onderstaande tabel zie je voor twaalf vluchten het cargogewicht (in duizenden kilogram) en het gemeten kerosineverbruik tijdens de cruise-fase (in kilogram per minuut)
    {\par\bigskip
        \begin{center}
            \begin{tabular}{c|cccccccccc}
                \toprule
                    \textbf{Cargogewicht (in duizenden kg)} & $2$ & $4$ & $6$ & $8$ & $10$ & $12$ & $14$ & $16$ & $18$ & $20$ \\
                \midrule
                    \textbf{Kerosineverbruik (in kg/min)} & $29$ & $17$ & $61$ & $76$ & $49$ & $81$ & $124$ & $102$ & $103$ & $118$ \\
                \bottomrule
            \end{tabular}
        \end{center}\par\bigskip
    }
}

    \subquestion{2}{
        Als we een regressie-analyse willen uitvoeren, welke variabele zou dan de afhankelijke variabele $Y$ zijn en welke variabele zou de onafhankelijke variabele $X$ zijn?
    }
    \solution{
        Bij een regressie-analyse is de meest logische keuze om cargogewicht als onafhankelijke variabele $X$ te kiezen, en het kerosineverbruik als afhankelijke variabele $Y$. \rubric{1}
        Een hoger kerosineverbruik verklaart niet het cargogewicht, terwijl dat andersom wellicht wel het geval kan zijn.\rubric{1}
    }

    \subquestion{5}{
        Teken het spreidingsdiagram op basis van je antwoord op vraag (a).
        Geef hierbij ook duidelijk aan welke as welke variabele beschrijft.
    }
    \solution{
        Het spreidingsdiagram is hieronder weergegeven, waarbij het cargogewicht (in kg) op de $x$-as is uitgezet en het kerosineverbruik (in kg/min) op de $y$-as.\rubric{1}
        \begin{center}
            \resizebox{0.9\textwidth}{!}{
                \includegraphics{exam20260519_regression.png} 
            }
        \end{center}\rubric{4}
    }

    \subquestion{8}{
        Bereken Pearson's correlatieco\"effici\"ent $r(x,y)$ op basis van de data in de bovenstaande tabel.
        Wat zegt $r(x, y)$ over de relatie tussen het cargogewicht en het kerosineverbruik?
    }
    \solution{
        We beginnen met het uitrekenen van Pearson's correlatieco\"effici\"ent 
        \begin{align*}
            r(x,y)  &= \frac{ \overline{x \cdot y} - \overline{x} \cdot \overline{y} }{ \sqrt{ (\overline{x^2} - \overline{x}^2 ) \cdot (\overline{y^2} - \overline{y}^2) }}
        \end{align*}
        Hiervoor hebben we dus de waardes van $\overline{x}$, $\overline{y}$, $\overline{xy}$, $\overline{x^2}$ en $\overline{y^2}$ nodig.
        Deze bepalen we aan de hand van de volgende rekentabel:
        \begin{center}
            \begin{tabular}{ccccc}
                \toprule
                    $x$ & $y$ & $xy$ & $x^2$ & $y^2$ \\
                \midrule
    		        $2$ & $29$ & $58$ & $4$ & $841$ \\
                    $4$ & $17$ & $68$ & $16$ & $289$ \\
                    $6$ & $61$ & $366$ & $36$ & $3721$ \\
                    $8$ & $76$ & $608$ & $64$ & $5776$ \\
                    $10$ & $49$ & $490$ & $100$ & $2401$ \\
                    $12$ & $81$ & $972$ & $144$ & $6561$ \\
                    $14$ & $124$ & $1736$ & $196$ & $15376$ \\
                    $16$ & $102$ & $1632$ & $256$ & $10404$ \\
                    $18$ & $103$ & $1854$ & $324$ & $10609$ \\
                    $20$ & $118$ & $2360$ & $400$ & $13924$ \\
                \midrule
                    $\overline{x} = 11$ & $\overline{y} = 76$ & $\overline{xy} = 1014.4$ & $\overline{x^2} = 154$ & $\overline{y^2} = 6990.2$ \rubric{4} \\
                \bottomrule
            \end{tabular}
        \end{center}
        \par\medskip
        De correlatieco\"effici\"ent van Pearson is dus gelijk aan
        \begin{align*}
            r(x,y)  &= \frac{ \overline{x \cdot y} - \overline{x} \cdot \overline{y} }{ \sqrt{ (\overline{x^2} - \overline{x}^2) \cdot (\overline{y^2} - \overline{y}^2) } }\\
                    &= \frac{ 1014.4 - 11 \cdot 76 }{ \sqrt{ (154 - 11^2) \cdot (6990.2 - 76^2) } } \\
                    &= \frac{178.4}{200.1714} \\
                    &\approx 0.8912. \rubric{2}
        \end{align*}
        De correlatieco\"effici\"ent $r \approx 0.8912$ ligt ver van $0$ af en is positief.
        Dit houdt in dat er een sterke positieve (stijgende) trend is. \rubric{1}
        Een hoger gewicht van de lading correleert sterk met een hoger kerosineverbruik. \rubric{1}
    }

    \subquestion{6}{
        Bereken de regressielijn $Y = a + b \cdot X$ door berekening van de co\"effici\"enten $a$ en $b$.
        Bepaal aan de hand van de regressielijn een statistisch verantwoorde voorspelling voor het kerosineverbruik van een transportvliegtuig met een lading van 5500 kilogram.
    }
    \solution{
        Om de co\"effici\"enten $a$ en $b$ van de regressielijn te berekenen, gaan we de rekentabel van vraag (a) hergebruiken.
        Er volgt namelijk dat
        \begin{align*}
            b &= \frac{\overline{xy} - \overline{x} \cdot \overline{y}}{\overline{x^2} - (\overline{x})^2} \\
              &= \frac{1014.4 - 11 \cdot 76}{154 - (11)^2} \\
              &= \frac{178.4}{33} \approx 5.4061, \rubric{2}\\
            a &= \overline{y} - b \cdot \overline{x} \\
              &= 76 - 5.4061 \cdot 11 \\
              &\approx 16.5333. \rubric{2}
        \end{align*}
        De formule van de regressielijn behorende bij deze steekproef is dus gelijk aan $Y = 16.5333 + 5.4061 \cdot X$. \rubric{1}
        Een statistisch verantwoorde voorspelling voor het kerosineverbruik van een transportvliegtuig met een lading van 5500 kilogram vinden we door $X = 5.5$ in te vullen:
        Dit geeft een waarde van $y_0 = 16.5333 + 5.4061 \cdot 5.5 \approx 46.2667$, oftewel net iets meer dan $46$ kilogram per minuut. \rubric{1} 
    }

    \subquestion{9}{
        Bereken een \SI{95}{\percent}-voorspellingsinterval voor het kerosineverbruik van een transportvliegtuig met een laadgewicht van 5500 kilogram.
        Wat kun je concluderen aan de hand van dit interval over de betrouwbaarheid van de voorspelling die je in vraag (d) hebt gemaakt?
    }
    \solution{
        In opdracht (d) hebben we een puntschatting van $y_0 = 16.5333 + 5.4061 \cdot 5.5 \approx 46.2667$ bepaald.
        Daarnaast kunnen we de standaardafwijking $\sigma$ van de storingsterm $\varepsilon$ schatten:
        \begin{align*}
            s_{\varepsilon} &= \sqrt{ \frac{n}{n-2} \cdot \left( \overline{y^2} - a \cdot \overline{y} - b \cdot \overline{xy} \right) } \\ 
                            &= \sqrt{ \frac{10}{8} \cdot \left( 6990.2 - 16.5333 \cdot 76 - 5.4061 \cdot 1014.4 \right) } \\ 
                            &\approx 17.6691. \rubric{2}
        \end{align*}

        Vervolgens kunnen we een puntschatting berekenen van de standaardafwijking van $Y$ voor gegeven $X = x_0$:
        \begin{align*}
            s_{f} &= s_{\varepsilon} \cdot \sqrt{ 1 + \frac{1}{n} \cdot \left( 1 + \frac{(x_0 - \overline{x})^2}{\overline{x^2} - \overline{x}^2} \right) } \\
                    &= 17.6691 \cdot \sqrt{ 1 + \frac{1}{10} \cdot \left( 1 + \frac{(5.5 - 11)^2}{154 - 11^2} \right) } \\
                    &\approx 19.2882. \rubric{2}
        \end{align*}

        Omdat we de standaardafwijkingen geschat hebben en de storingstermen normaal verdeeld zijn, moeten we werken met de $t$-verdeling met $df = n - 2 = 8$ vrijheidsgraden.
        De $t$-waarde die hoort bij een betrouwbaarheidsniveau $\alpha = 0.05$ is gelijk aan
        \[
            t = \invt(\opp = 1 - \alpha / 2, \df = n - 2) = \invt(\opp = 0.975, \df = 8) \approx 2.306.\rubric{1}
        \]
        Het \SI{95}{\percent}-betrouwbaarheidsinterval voor de gemiddelde $Y$ voor gegeven $X = x_0$ kan dus worden beschreven door
        \begin{align*}
            &[y_0 - t \cdot s_{f}, y_0 - t \cdot s_{f}] \\ 
            &= [46.2667 - 2.306 \cdot 19.2882; 46.2667 + 2.306 \cdot 19.2882] \\ 
            &\approx [1.7879, 90.7454]. \rubric{2}
        \end{align*}
        
        Met \SI{95}{\percent} betrouwbaarheid zal voor een gegeven $X = 5.5$ een bijbehorende uitkomst van $Y$ liggen tussen $1$ en $91$. 
        Dit interval is erg breed, wat betekent dat de voorspelling van $Y\approx 46.2667$ voor gegeven $X = 5.5$ niet per se nauwkeurig hoeft te zijn. \rubric{2}
    }
\end{question}